\section*{Chapitre \ref{ch:rs:trade-flow-features} --- Caractéristiques du flux de transactions}

\begin{solution}{exo:rs:trade-flow-features:1}
Le prix milieu vaut 20.005. \omterm{def:rs:trade-flow-features:sign}{Règle du tick}~: non classée (pas de prix antérieur), $+1$, $+1$ (un tick nul conserve le dernier signe), $-1$, $-1$. \omterm{def:rs:trade-flow-features:sign}{Algorithme de Lee--Ready}~: chaque
transaction est située hors du prix milieu, c'est donc la \omterm{def:rs:trade-flow-features:sign}{règle de la cotation}~: $-1$, $+1$, $+1$, $-1$, $-1$. Les deux ne diffèrent que sur la première transaction.
\end{solution}

\begin{solution}{exo:rs:trade-flow-features:2}
$\Phi(0.8) = 0.788$~: 78.8~\% du volume de la barre est classé en achats, quelles qu'aient été en réalité les transactions qu'elle contient.
\end{solution}

\begin{solution}{exo:rs:trade-flow-features:3}
Exposant de décroissance $\gamma = \alpha - 1 = 0.4$~; exposant de Hurst $H = (3 - \alpha)/2 = 0.8$.
\end{solution}

\begin{solution}{exo:rs:trade-flow-features:4}
Les déséquilibres $|V^B - V^S| = |2V^B - 1\,000|$ valent 400, 200, 800, 0 et 600, leur somme vaut 2\,000, d'où le \omterm{def:rs:trade-flow-features:vpin}{VPIN} $= 2\,000/(5 \times 1\,000) = 0.4$.
\end{solution}

\begin{solution}{exo:rs:trade-flow-features:5}
$30 \times 0.10 = 3.0$ ticks. Le volume signé des transactions ne voit que les exécutions~; les cotations bougent aussi lorsque des ordres sont annulés ou ajoutés, ce que
l'OFI retient et le volume signé non. Sur le tape simulé, le $R^2$ vaut 14~\% contre 53~\% pour l'OFI, si bien que la pente est estimée
à partir d'une relation qui laisse inexpliquée la plus grande part de la variation du prix.
\end{solution}

\begin{solution}{exo:rs:trade-flow-features:6}
Prenons un processus de Poisson d'intensité $\lambda_1$ pendant la moitié du temps et $\lambda_2 > \lambda_1$ pendant l'autre moitié, chaque régime étant long. Aucun événement ne
cause un autre, mais un événement a davantage de chances de tomber dans le régime actif, et davantage d'événements le suivent alors vite~: les comptages sont sur-dispersés
par rapport à un seul processus de Poisson. Un processus de Hawkes à intensité de fond constante n'a qu'une seule façon de produire ce regroupement, l'excitation, si bien que
l'ajustement place de la masse dans le noyau, et plus les régimes diffèrent, plus le ratio de branchement est élevé. Sur le tape du chapitre, l'ajustement donne 0.65
contre 0.4 imposé, valeur que le tape plat retrouve (0.403).
\end{solution}

\begin{solution}{exo:rs:trade-flow-features:7}
\texttt{rs\_tradeflow.stale\_quotes}~: 99.7~\% avec des cotations en retard de 0.05 seconde, 97.9~\% à 0.5 seconde et 93.8~\% à 2 secondes. La perte vaut environ
cinq points par seconde de retard au départ, et elle ralentit (6.2 points à deux secondes)~: plus la cotation est ancienne, plus il est probable que le prix milieu ait bougé
au-delà du prix de la transaction, mais une cotation qui a bougé une fois peut revenir en arrière.
\end{solution}

\begin{solution}{exo:rs:trade-flow-features:8}
C'est une affirmation choisie après l'événement. Une alerte n'est utile qu'avec son taux de fausses alarmes~: à quelle fréquence le \omterm{def:rs:trade-flow-features:vpin}{VPIN} a atteint de tels niveaux sans krach,
et si les valeurs élevées étaient suivies, en moyenne, d'une forte volatilité (Andersen et Bondarenko ont constaté qu'elles étaient un mauvais prédicteur de celle-ci). Une
lecture maximale de l'année n'est elle aussi qu'un tirage parmi les heures de l'année~; une heure quelconque avant un événement est le maximum de son voisinage. Et le
niveau dépend de choix (taille des paniers, fenêtre, méthode de signature des transactions) qui doivent être fixés avant l'événement pour que l'affirmation vaille.
\end{solution}

\begin{solution}{pb:rs:trade-flow-features:1}
\begin{enumerate}
\item 15~\% avant l'épisode, 45~\% pendant celui-ci~: la part des informés est presque multipliée par trois.
\item $-0.36$ ticks avant, $+0.96$ pendant (le prix milieu dix secondes après la transaction, moins le prix de la transaction, dans le sens de la transaction).
\item Les agresseurs paient la demi-fourchette, soit un demi-tick, et la plupart d'entre eux sont des traders de bruit dont les ordres ne portent aucune information~: le
côté passif encaisse la fourchette.
\item Oui~: les agresseurs gagnent et les fournisseurs de liquidité perdent en moyenne, ce qui est la définition du flux toxique.
\item Paniers de 2\,649 actions, soit $1/400$ du volume de la séance, donc 400 paniers, et une fenêtre de 20 paniers.
\item 0.31 avant, 0.35 pendant~: une hausse de 13~\%.
\item À la minute 119.7, soit 20 minutes après la fin de l'épisode.
\item 16 secondes après le début de l'épisode~; le \omterm{def:rs:trade-flow-features:vpin}{VPIN} dépasse le même seuil pendant 6.0~\% de la séance en dehors de l'épisode.
\item Les traders informés achètent lorsque le prix efficient est au-dessus du prix milieu et vendent lorsqu'il est en dessous~; comme ce prix oscille dans les deux sens, leur flux
sur un panier (27 secondes en moyenne) n'est pas unilatéral. Les séquences unilatérales proviennent des \omterm{def:rs:trade-flow-features:acf}{métaordres} des traders de bruit, que l'épisode remplace
par un flux informé plutôt que de s'y ajouter.
\item $-0.38$ avant, $0.84$ pendant~; il atteint son maximum à 5\,740 secondes, à l'intérieur de l'épisode.
\item À 5\,440 secondes, soit 40 secondes après le début~: chaque valeur attend dix secondes pour son mark-out, et la moyenne glissante sur une minute a besoin
de transactions toxiques pour se remplir.
\item 8.4~\%.
\item Pour étudier la toxicité a posteriori, le mark-out~: c'est la toxicité par définition. Pour protéger un teneur de marché en temps réel, les mark-outs
de ses propres exécutions récentes à un horizon court, avec les caractéristiques du carnet d'ordres du chapitre~8~; le \omterm{def:rs:trade-flow-features:vpin}{VPIN} n'a pas montré ici qu'il aiderait.
\item Des traders informés disposant d'un signal qui pointe dans un sens plus longtemps qu'un panier, un unique gros ordre informé fractionné sur plusieurs minutes, ou une nouvelle
qui fait bouger la valeur en un seul saut. Le \omterm{def:rs:trade-flow-features:vpin}{VPIN} mesure ce cas-là, et non celui-ci.
\item 0.65 sur ce tape~; 0.403 sur le tape plat, contre les 0.4 imposés.
\item 0.10 tick par 100 actions d'achats nets, avec un $R^2$ de 14~\%.
\item Le \omterm{def:rs:trade-flow-features:vpin}{VPIN} franchit son seuil en premier, au bout de 16 secondes, avec 6.0~\% de fausses alarmes~; le mark-out, au bout de 40 secondes, avec 8.4~\%. Aucun des deux n'est propre, et
le franchissement précoce du \omterm{def:rs:trade-flow-features:vpin}{VPIN} ne s'accompagne pas d'une hausse de son niveau.
\item \textbf{Résultat nommé.} Lorsque la part informée du flux passe de 15~\% à 45~\% et que le mark-out des agresseurs passe de $-0.36$ à
$+0.96$ ticks, le \omterm{def:rs:trade-flow-features:vpin}{VPIN} ne progresse que de 13~\%, de 0.31 à 0.35, et atteint son maximum 20 minutes après la fin de l'épisode.
\item Une mesure de toxicité doit être testée contre la toxicité elle-même (les mark-outs), sur des données où elle est connue, avant qu'on lui fasse confiance là où elle
ne l'est pas~; une histoire plausible reliant la mesure au trading informé n'est pas un test.
\item Le flux toxique est un flux d'ordres dont les agresseurs sont informés, de sorte que les fournisseurs de liquidité qui négocient contre lui perdent en moyenne.
\end{enumerate}
\end{solution}

\begin{solution}{iq:rs:trade-flow-features:1}
Avec des cotations, comparer le prix de la transaction au prix milieu en vigueur, en prenant les cotations à l'horodatage de la transaction et en vérifiant comment les deux horloges
s'alignent (Lee et Ready, 1991)~; se rabattre sur la \omterm{def:rs:trade-flow-features:sign}{règle du tick} au prix milieu. Sans cotations, la \omterm{def:rs:trade-flow-features:sign}{règle du tick}. Mesurer l'erreur là où la vérité est
connue (un flux de données avec un indicateur d'agresseur, ou un simulateur)~: sur le tape du chapitre, une cotation périmée d'une seconde fait perdre 3.7 points de précision.
\end{solution}

\begin{solution}{iq:rs:trade-flow-features:2}
Les \omterm{def:rs:trade-flow-features:acf}{métaordres} sont découpés en de nombreux ordres enfants de même signe, si bien que les signes persistent (Lillo, Mike et Farmer, 2005). Les prix restent presque
imprévisibles, car la liquidité s'ajuste~: à mesure qu'une série d'achats se poursuit, l'impact de chaque achat supplémentaire baisse, et les vendeurs reconstituent le cours vendeur, si bien que le signe prévisible
ne se transforme pas en un mouvement de prix prévisible de même ampleur.
\end{solution}

\begin{solution}{iq:rs:trade-flow-features:3}
Un flux est toxique lorsque ceux qui négocient contre lui perdent en moyenne. Le mesurer par les mark-outs~: pour chacune des exécutions du carnet, le prix milieu quelques horizons
plus tard, moins le prix d'exécution, du côté du carnet, par client, plateforme, moment de la journée et taille~; un mark-out qui devient négatif rapidement est un signe de
flux toxique.
\end{solution}

\begin{solution}{iq:rs:trade-flow-features:4}
Une intensité de fond qui varie (moment de la journée, nouvelles, régimes de volatilité) produit des grappes qu'un ajustement à intensité de fond constante attribue à l'excitation. Ajuster
avec une intensité de fond qui suit le moment de la journée ou un niveau d'activité mesuré, comparer sur des données hors échantillon, et vérifier sur des données simulées avec un
ratio de branchement connu~: le tape du chapitre donne 0.65, sa version plate 0.403 contre 0.4.
\end{solution}

\begin{solution}{iq:rs:trade-flow-features:5}
Découper les transactions en paniers de volume égal, leur affecter un signe (ou classer les barres par BVC), et moyenner le déséquilibre absolu achat--vente sur les $n$
derniers paniers, divisé par la taille du panier~; l'affirmation est que le flux unilatéral est toxique. La critique~: ce fut un mauvais prédicteur de la volatilité à court terme
(Andersen et Bondarenko, 2014), son niveau dépend de choix arbitraires, et un flux unilatéral n'est pas forcément informé, ni un flux informé
forcément unilatéral.
\end{solution}

\begin{solution}{iq:rs:trade-flow-features:6}
Si les tailles des \omterm{def:rs:trade-flow-features:acf}{métaordres} ont un exposant de queue $\alpha$, l'autocorrélation du signe décroît en $\tau^{-(\alpha-1)}$ (Lillo, Mike et Farmer, 2005), et une
série dont l'autocorrélation décroît en $\tau^{-\gamma}$ avec $0 < \gamma < 1$ a des sommes partielles dont la variance croît comme $n^{2 - \gamma}$,
donc $H = 1 - \gamma/2 = (3 - \alpha)/2$. Test~: agréger les exécutions en ordres, calculer la variance des sommes par blocs des signes à travers les échelles
et effectuer la régression en logarithmes. Le tape du chapitre, avec $\alpha = 1.5$, prédit 0.75 et donne 0.715.
\end{solution}
